\subsubsection{从自对偶压缩到一变量 \texorpdfstring{$\Xi$}{Xi}：严格函数方程与临界线实值性}\label{subsec:xi-from-self-dual}
本节把附录中“自对偶 + 完成化 + 单位圆实值化”的三件套压缩为一个一变量 \(\Xi\)-对象，并将函数方程理解为完成化投影系统上的自同构在截面上的自然性约束；相应地，“临界线”不再被视作 \(s\)-平面上的对称轴，而被解释为参数侧的 unitary slice（\(|u|=1\)）这一协议可见切片。

\paragraph{自对偶条件}
设 \(\Delta(z,u)\) 是二变量多项式/解析函数（在本文的核口径下典型为
\(\Delta(z,u)=\det(I-zB(u))\)，见 \S\ref{subsec:zeta-online-kernel}），并满足二变量自对偶恒等式
\[
\boxed{\ \Delta(z,u)=\Delta(uz,1/u)\ }.
\]
对同步核加权情形，该恒等式在附录 \ref{app:kernel-self-dual} 给出；并由此导出完成化偶性与系数回文约束（例如推论 \ref{cor:delta-coeff-palindrome}）。

\paragraph{新的一变量对象：任意底 \texorpdfstring{$L>1$}{L>1}}
取任意常数底 \(L>1\)，定义
\[
z_L(s):=L^{-s},\qquad u_L(s):=L^{2s-1},
\]
并定义一变量函数
\[
\boxed{\ \Xi_{\Delta,L}(s):=\Delta\!\bigl(z_L(s),u_L(s)\bigr)=\Delta\!\bigl(L^{-s},L^{2s-1}\bigr)\ }.
\]
当取 \(L=\lambda=\rho(B(1))\) 时，这与附录 \ref{app:kernel-functional-equation} 的截面 \(\Xi(s)\) 一致（命题 \ref{prop:kernel-xi-functional-equation}）。

\begin{proposition}[严格函数方程：投影系统自同构的截面自然性]\label{prop:xi-delta-L-functional-equation}
在上述自对偶条件下，对任意 \(L>1\) 都有严格恒等式
\[
\boxed{\ \Xi_{\Delta,L}(s)=\Xi_{\Delta,L}(1-s)\ }.
\]
\end{proposition}
\begin{proof}
令 \(\tau:(z,u)\mapsto(uz,u^{-1})\)。自对偶恒等式 \(\Delta(z,u)=\Delta(uz,1/u)\) 等价于 \(\Delta=\Delta\circ\tau\)。
对任意 \(L>1\) 直接计算得
\[
u_L(1-s)=u_L(s)^{-1},\qquad z_L(1-s)=u_L(s)\,z_L(s),
\]
从而 \((z_L(1-s),u_L(1-s))=\tau(z_L(s),u_L(s))\)。
因此
\[
\Xi_{\Delta,L}(s)
=\Delta\!\bigl(z_L(s),u_L(s)\bigr)
=\Delta\!\Bigl(\tau\bigl(z_L(s),u_L(s)\bigr)\Bigr)
=\Delta\!\bigl(z_L(1-s),u_L(1-s)\bigr)
=\Xi_{\Delta,L}(1-s).
\]
\end{proof}

\begin{remark}[投影系统自同构：函数方程的对象]\label{rem:xi-projection-system-automorphism}
自对偶恒等式 \(\Delta(z,u)=\Delta(uz,1/u)\) 可视为二变量接口上的一个 involution：
定义
\[
\tau:\ (z,u)\longmapsto (uz,u^{-1}),
\qquad \tau^2=\mathrm{id}.
\]
则自对偶条件等价于 \(\Delta=\Delta\circ\tau\)，即 \(\Delta\) 在完成化投影系统的自同构 \(\tau\) 下不变。
\par
对任意 \(L>1\)，完成化截面映射 \(s\mapsto (z_L(s),u_L(s))\) 与 \(s\mapsto 1-s\) 在该接口上严格对齐：
\[
(z_L(1-s),u_L(1-s))
=(L^{s-1},L^{1-2s})
\;=\;
\bigl(u_L(s)\,z_L(s),\,u_L(s)^{-1}\bigr)
=\tau\bigl(z_L(s),u_L(s)\bigr).
\]
因此，命题 \ref{prop:xi-delta-L-functional-equation} 的函数方程并非“关于某条几何轴的对称”，而是 \(\tau\)-不变性在完成化截面上的拉回。该口径也解释了为何 unitary-slice 核验协议自然锁定 \(|u|=1\) 的可见域：这是一条与 \(\tau\) 相容、并支撑 Hardy 型实值化与谱穿越审计的参数切片。
\end{remark}

\paragraph{完成化与“临界线实值性”}
写
\[
\Delta(z,u)=\sum_{j=0}^{N}c_j(u)\,z^j.
\]
若系数满足完成化回文结构 \(c_j(u)=u^j c_j(1/u)\)（推论 \ref{cor:delta-coeff-palindrome}），则可定义完成化系数
\[
d_j(u):=u^{-j/2}c_j(u).
\]
附录推论 \ref{cor:completion-real-line} 表明：当 \(|u|=1\) 时，\(d_j(u)\in\RR\)（完成化后“单位圆上系数为实”）。

\begin{proposition}[临界线实值性：完成化 \(\Rightarrow\) Hardy 型实值化]\label{prop:xi-delta-L-critical-line-real}
在自对偶条件且满足推论 \ref{cor:completion-real-line} 的条件下，对任意 \(L>1\) 与任意实数 \(t\) 有
\[
\boxed{\ \Xi_{\Delta,L}\!\left(\tfrac12+it\right)\in\RR\ }.
\]
\end{proposition}
\begin{proof}
取 \(s=\tfrac12+it\)，则 \(u_L(s)=L^{2it}\) 满足 \(|u_L(s)|=1\)，并且
\[
\Xi_{\Delta,L}(s)
=\Delta(L^{-s},u_L(s))
=\sum_{j=0}^{N}c_j(u_L(s))\,L^{-js}
=\sum_{j=0}^{N}d_j(u_L(s))\,\bigl(L^{-1/2}\bigr)^{j}\,e^{-ijt\log L}.
\]
由 \(|u_L(s)|=1\Rightarrow d_j(u_L(s))\in\RR\)，上述表达是一个关于实变量 \(e^{-it\log L}\) 的“实系数”三角多项式在 \(t\) 处的取值，因此其值为实数。等价地，可把附录中 Hardy 实值化 \(Z(t)\) 的证明（命题 \ref{prop:kernel-hardy-real-even}）中的 \(\log\lambda\) 替换为 \(\log L\)。
\end{proof}

\begin{remark}[临界线 \texorpdfstring{$\Leftrightarrow$}{<->} 单位圆切片]\label{rem:xi-critical-line-unit-circle}
由 \(u_L(s)=L^{2s-1}\) 可见：直线 \(\Re(s)=\tfrac12\) 恰对应 \(|u_L(s)|=1\)。
因此在本文的可核验语义中，“临界线”并非额外几何公理，更不是 \(s\)-平面上的对称轴；它是完成化投影系统在参数侧的 unitary slice，并在协议层被固定为可见域 \(X_{\mathrm{vis}}=\mathsf{Phase}\) 的唯一连续切片（定义 \ref{def:xi-visible-read-null}）。
结合注 \ref{rem:xi-projection-system-automorphism}，函数方程的结构对象是投影系统自同构 \(\tau\)，而 unitary slice 则是该自同构在可见域投影下保持稳定的审计切片：在该切片上完成化系数可取实，从而支持 Hardy 型实值化与谱穿越接口。
\end{remark}

\begin{proposition}[unitary slice 的行列式重写与谱穿越判据]\label{prop:xi-unitary-slice-spectral-crossing}
在自对偶条件下，进一步假设
\(
\Delta(z,u)=\det(I-zB(u))
\)
来自某个有限维（可审计）矩阵族 \(B(u)\)（例如同步核加权情形，见 \S\ref{subsec:zeta-online-kernel} 与附录 \ref{app:kernel-self-dual}）。
固定 \(L>1\)，令 \(\theta:=2t\log L\)。
则对任意 \(t\in\RR\) 有严格恒等式
\[
\boxed{\;
\Xi_{\Delta,L}\!\left(\tfrac12+it\right)
=
\det\!\Bigl(I-U_L(\theta)\Bigr),\qquad
U_L(\theta):=L^{-1/2}e^{-i\theta/2}\,B(e^{i\theta})\; }.
\]
因此临界线零点满足等价判据
\[
\boxed{\;
\Xi_{\Delta,L}\!\left(\tfrac12+it\right)=0
\quad\Longleftrightarrow\quad
1\in\mathrm{Spec}\bigl(U_L(\theta)\bigr),\ \ \theta=2t\log L\; }.
\]
\end{proposition}
\begin{proof}
取 \(s=\tfrac12+it\)。则
\(
u_L(s)=L^{2it}=e^{i\theta}
\)
其中 \(\theta=2t\log L\)，并且
\(
z_L(s)=L^{-s}=L^{-1/2}e^{-it\log L}=L^{-1/2}e^{-i\theta/2}
\)。
由 \(\Xi_{\Delta,L}(s)=\Delta(z_L(s),u_L(s))=\det(I-z_L(s)B(u_L(s)))\) 得到所示重写。
零点判据由行列式为零当且仅当 \(1\) 为 \(U_L(\theta)\) 的特征值立得。证毕。
\end{proof}

\begin{remark}[可审计的谱穿越口径：从多项式根到特征值轨道]\label{rem:xi-unitary-slice-operator-audit}
命题 \ref{prop:xi-unitary-slice-spectral-crossing} 把 Hardy 截面上的零点检测从
“一元多项式 \(P_L(S)\) 的根”（命题 \ref{prop:xi-rh-interval-criterion}）
等价转写为
“unitary 参数 \(\theta\) 下归一化核 \(U_L(\theta)\) 的谱是否穿过 \(1\)”。
这一口径的意义在于：它把“临界线几何”从 \(s\)-平面的隐式条件，转化为 \(|u|=1\) 上的显式算子族谱问题，从而允许直接以数值线性代数在 \(\theta\)-轴上进行独立审计。
\par
同步核加权例子中，脚本 \path{scripts/exp_sync_kernel_unitary_slice_spectral_crossing.py} 在离散网格上同时计算
\(\det(I-U_L(\theta))\) 与 \(\min_{\lambda\in\mathrm{Spec}(U_L(\theta))}|\lambda-1|\)，并输出为可复现表格：
\input{sections/generated/tab_sync_kernel_unitary_slice_spectral_crossing}
\end{remark}

\paragraph{自对偶截面上的完成化常数：\texorpdfstring{$w$}{w} 固定，\texorpdfstring{$\Xi$}{Xi} 退化为余弦多项式}
为避免与复变量 \(s\) 混淆，本段把 completion 的迹坐标记为 \(S\)。
若自对偶 \(\Delta(z,u)\) 还可写成某个完成化行列式
\(\widehat\Delta(w,S)\) 的形式
\[
\Delta(z,u)=\widehat\Delta\!\left(z\sqrt{u},\ \sqrt{u}+\frac{1}{\sqrt{u}}\right)
\]
其中 \(\widehat\Delta\) 在自对偶核上由不变量子环强制存在；例如同步核加权情形见命题 \ref{prop:sync-kernel-weighted-hat-delta}。
则沿自对偶 Dirichlet 截面 \(z=L^{-s},\ u=L^{2s-1}\) 有
\[
w:=z\sqrt{u}=L^{-s}\cdot L^{s-1/2}=L^{-1/2}\quad(\text{常数}),
\qquad
S_L(s):=\sqrt{u}+\frac{1}{\sqrt{u}}=L^{s-1/2}+L^{1/2-s},
\]
从而得到一维压缩恒等式
\[
\boxed{\ \Xi_{\Delta,L}(s)=\widehat\Delta\!\bigl(L^{-1/2},\,S_L(s)\bigr)\ }.
\]
特别地在临界线 \(s=\tfrac12+it\) 上，
\[
S_L\!\left(\tfrac12+it\right)=L^{it}+L^{-it}=2\cos(t\log L),
\]
因此 Hardy 型截面 \(\Xi_{\Delta,L}(\tfrac12+it)\) 是一个“余弦的实多项式”。

\begin{proposition}[零点临界线版的纯代数判据：\texorpdfstring{$\Xi$}{Xi}-RH \(\Longleftrightarrow\) 根区间]\label{prop:xi-rh-interval-criterion}
在上述完成化表示下，令
\[
P_L(S):=\widehat\Delta\!\bigl(L^{-1/2},S\bigr)\in\RR[S].
\]
则下述两条陈述等价：
\begin{enumerate}
  \item \(\Xi_{\Delta,L}\) 的全部零点都满足 \(\Re(s)=\tfrac12\)；
  \item \(P_L(S)\) 的全部根都落在实区间 \([-2,2]\) 内。
\end{enumerate}
\end{proposition}
\begin{proof}
令 \(y:=L^{s-1/2}\in\CC^\times\)。则 \(S_L(s)=y+y^{-1}\)，从而
\[
\Xi_{\Delta,L}(s)=\widehat\Delta\!\bigl(L^{-1/2},\,y+y^{-1}\bigr)=P_L(y+y^{-1}).
\]
因此 \(\Xi_{\Delta,L}(s)=0\) 当且仅当 \(S_0:=y+y^{-1}\) 为 \(P_L\) 的某个根。
\par
另一方面，
\[
|y|=|L^{s-1/2}|=L^{\Re(s)-1/2},
\]
故 \(\Re(s)=\tfrac12\) 当且仅当 \(|y|=1\)。
而 \(|y|=1\) 当且仅当 \(y=e^{i\theta}\)（某个 \(\theta\in\RR\)），此时
\(
y+y^{-1}=2\cos\theta\in[-2,2]
\) 为实数。
反之，若 \(S_0\in[-2,2]\) 为实数，则方程 \(y+y^{-1}=S_0\) 的解满足 \(|y|=1\)。
因此 “所有零点都满足 \(\Re(s)=\tfrac12\)” 等价于 “\(P_L\) 的所有根都在 \([-2,2]\)”。
\end{proof}

\begin{corollary}[单位圆根判据：\texorpdfstring{$\Xi$}{Xi}-RH \(\Longleftrightarrow\) 自倒数多项式的单位圆根分布]\label{cor:xi-rh-unit-circle-reciprocal}
在命题 \ref{prop:xi-rh-interval-criterion} 的记号下，令 \(d:=\deg P_L\) 并定义
\[
\Phi_L(y):=y^{d}P_L(y+y^{-1})\in\RR[y].
\]
则：
\begin{enumerate}
  \item \(\Phi_L\) 为自倒数（reciprocal/palindromic）多项式，即
  \[
  \boxed{\ y^{2d}\Phi_L(1/y)=\Phi_L(y)\ }.
  \]
  \item 下列陈述等价：
  \begin{enumerate}
    \item[(i)] \(\Xi_{\Delta,L}\) 的全部零点都满足 \(\Re(s)=\tfrac12\)；
    \item[(ii)] \(\Phi_L(y)\) 的全部复根都满足 \(|y|=1\)；
    \item[(iii)] \(P_L(S)\) 的全部根都落在实区间 \([-2,2]\) 内。
  \end{enumerate}
\end{enumerate}
\end{corollary}
\begin{proof}
由定义 \(\Phi_L(y)=y^{d}P_L(y+y^{-1})\)，把 \(y\mapsto 1/y\) 代入得
\(
\Phi_L(1/y)=y^{-d}P_L(y+y^{-1})
\)，
故 \(y^{2d}\Phi_L(1/y)=y^{d}P_L(y+y^{-1})=\Phi_L(y)\)，从而 \(\Phi_L\) 自倒数。
\par
另一方面，\(\Phi_L(y)=0\) 当且仅当 \(P_L(S_0)=0\) 且 \(S_0=y+y^{-1}\)。
若 \(S_0\in[-2,2]\)，则 \(y\) 必可写为 \(y=e^{i\theta}\)（\(\theta\in\RR\)），从而 \(|y|=1\)；反之若 \(|y|=1\)，则 \(S_0=y+y^{-1}\in[-2,2]\)。
因此 “\(P_L\) 的根全在 \([-2,2]\)” 等价于 “\(\Phi_L\) 的根全在单位圆”，再与命题 \ref{prop:xi-rh-interval-criterion} 合并即得 (i)\(\Leftrightarrow\)(ii)\(\Leftrightarrow\)(iii)。
\end{proof}

\begin{proposition}[Cayley 超双曲化：单位圆根判据的实轴等价]\label{prop:xi-rh-cayley-hyperbolicity}
在推论 \ref{cor:xi-rh-unit-circle-reciprocal} 的记号下，令 \(D:=\deg\Phi_L=2d\)。
定义 Cayley 变换
\begin{equation}\label{eq:xi-cayley-y-from-x}
\boxed{\;
y=\mathcal{C}(x):=\frac{x-\mathrm{i}}{x+\mathrm{i}},
\qquad
x=\mathcal{C}^{-1}(y)=\mathrm{i}\,\frac{1+y}{1-y}\;}
\end{equation}
（它把 \(\overline{\RR}=\RR\cup\{\infty\}\) 双射到单位圆 \(\TT\)，并满足 \(\mathcal{C}(\infty)=1\)）。
并定义一元多项式
\begin{equation}\label{eq:xi-cayley-psi}
\boxed{\;
\Psi_L(x):=(x+\mathrm{i})^{D}\,\Phi_L\!\left(\frac{x-\mathrm{i}}{x+\mathrm{i}}\right)\; }.
\end{equation}
则下列条件等价：
\begin{enumerate}
  \item \(\Phi_L(y)\) 的全部复根都满足 \(|y|=1\)；
  \item \(\Psi_L(x)\) 的全部复根都落在 \(\RR\cup\{\infty\}\)（按射影意义计入 \(x=\infty\) 的根）。
\end{enumerate}
当 \(\Phi_L\in\RR[y]\) 为自倒数多项式时，\(\Psi_L\) 的系数经归一化可取为实数，因此上述条件可表述为 \(\Psi_L\) 的一元实\textbf{超双曲性}（全实根）。
\end{proposition}
\begin{proof}
由 \eqref{eq:xi-cayley-y-from-x} 可直接验证：对 \(x\in\RR\)，有 \(|\mathcal{C}(x)|=1\)；反之对 \(|y|=1\) 且 \(y\neq 1\)，有 \(x=\mathrm{i}(1+y)/(1-y)\in\RR\)，并且 \(y=1\) 对应 \(x=\infty\)。
因此 Cayley 变换在根集合层面给出双射
\(
\mathrm{Roots}(\Phi_L)\subseteq\TT
\Longleftrightarrow
\mathrm{Roots}(\Psi_L)\subseteq\overline{\RR}
\)：
若 \(\Phi_L(y_0)=0\) 且 \(y_0=\mathcal{C}(x_0)\)（\(x_0\in\RR\) 或 \(x_0=\infty\)），则清分母得 \(\Psi_L(x_0)=0\)（对 \(x_0=\infty\) 以射影根计入）；反之若 \(\Psi_L(x_0)=0\) 则 \(y_0=\mathcal{C}(x_0)\) 为 \(\Phi_L\) 的根且 \(|y_0|=1\)。
\par
若 \(\Phi_L\) 为自倒数且实系数，写 \(\Phi_L(y)=\sum_{k=0}^{D}a_k y^k\) 且 \(a_k=a_{D-k}\in\RR\)。
由 \eqref{eq:xi-cayley-psi} 展开得
\[
\Psi_L(x)=\sum_{k=0}^{D} a_k\,(x+\mathrm{i})^{D-k}(x-\mathrm{i})^{k}.
\]
利用对称性 \(a_k=a_{D-k}\) 成对合并 \(k\) 与 \(D-k\) 项即可把右端写成 \(\sum_k 2a_k\,\Re\bigl((x+\mathrm{i})^{D-k}(x-\mathrm{i})^{k}\bigr)\)，从而 \(\Psi_L\)（经整体常数归一化）具有实系数。证毕。
\end{proof}

\begin{proposition}[Cohn 型导数盘内性判别：单位圆根的等价证书]\label{prop:xi-rh-cohn-derivative}
在推论 \ref{cor:xi-rh-unit-circle-reciprocal} 的记号下，\(\Phi_L\in\RR[y]\) 为自倒数多项式且 \(\Phi_L(0)=\mathrm{lc}(P_L)\neq 0\)。
则下列陈述等价：
\begin{enumerate}
  \item \(\Phi_L\) 的全部复根满足 \(|y|=1\)；
  \item \(\Phi_L\) 为 self-inversive，并且其导数 \(\Phi_L'(y)\) 的全部复根满足 \(|y|\le 1\)。
\end{enumerate}
因此 \(\Xi\)-RH 条件（命题 \ref{prop:xi-rh-interval-criterion} / 推论 \ref{cor:xi-rh-unit-circle-reciprocal} 的任一等价形式）可在 \(|y|\le 1\) 的\textbf{导数盘内性}上给出。
\end{proposition}
\begin{proof}
自倒数与实系数蕴含
\(
\Phi_L(y)=y^{D}\overline{\Phi_L(1/\overline{y})}
\)
（其中 \(D=\deg\Phi_L\)），即 \(\Phi_L\) 为 self-inversive。
由 Cohn 定理：多项式的全部零点在单位圆上当且仅当其为 self-inversive 且其导数的全部零点落在闭单位圆盘内；参见 \cite{LalinSmyth2012UnimodularitySelfInversive}（其中回顾并推广了 Cohn 的原始判别）。
\end{proof}

\begin{remark}[Schur--Cohn 半正定证书接口：把“单位圆根/区间根”转写为可审计 LMI]\label{rem:xi-schur-cohn-sdp-interface}
令 \(p(y)=a_0+a_1y+\cdots+a_n y^n\in\CC[y]\)（\(a_0\neq 0\)）。
定义下三角 Toeplitz 矩阵
\[
X(p):=\begin{pmatrix}
a_0 & 0 & \cdots & 0\\
a_1 & a_0 & \ddots & \vdots\\
\vdots & \ddots & \ddots & 0\\
a_{n-1} & \cdots & a_1 & a_0
\end{pmatrix},
\qquad
Y(p):=\begin{pmatrix}
\overline{a_n} & 0 & \cdots & 0\\
\overline{a_{n-1}} & \overline{a_n} & \ddots & \vdots\\
\vdots & \ddots & \ddots & 0\\
\overline{a_1} & \cdots & \overline{a_{n-1}} & \overline{a_n}
\end{pmatrix},
\]
并令 Schur--Cohn Hermitian 形式矩阵
\[
H(p):=X(p)X(p)^\ast-Y(p)Y(p)^\ast\in\Mat_n(\CC)_{\mathrm{Herm}}.
\]
Schur--Cohn 理论给出：若 \(p\) 在单位圆上无零点，则 \(p\) 的全部零点落在开单位圆盘 \(|y|<1\) 当且仅当 \(H(p)\succ 0\)；并且 \(H(p)\) 的惯性指数给出单位圆内/外零点计数（标量情形的综述见 \cite{DymYoung1990SchurCohnMatrixPolynomials}）。
\par
结合命题 \ref{prop:xi-rh-cohn-derivative}，在典型非退化情形（例如 \(\Phi_L\) 的单位圆根均为单根）下，\(\Xi\)-RH 等价于 \(H(\Phi_L')\succ 0\)；
而在允许重根的边界情形，可将 \(|y|\le 1\) 的闭盘条件替换为严格版 \(|y|<1\)，或将 \(H(\Phi_L')\succ 0\) 放松为半正定 \(H(\Phi_L')\succeq 0\) 并配合边界根的有限审计。
\par
当 \(\Phi_L\)（或更一般的通道多项式 \(\Phi_{L,\rho}\)）的系数落在不变量子环 \(\QQ[u+u^{-1}]\) 中时，沿单位圆 \(u=e^{it}\) 可把系数视为 \(x=2\cos t\in[-2,2]\) 的多项式；于是 \(H(\Phi_{L,\rho}')\) 成为一个\textbf{一元多项式矩阵} \(H(x)\)，从而“全角度”核版 (G)RH 可压缩为区间半正定约束
\[
H(x)\succeq 0\qquad(x\in[-2,2]).
\]
该区间半正定性进一步可通过 SOS/半定规划给出可审计证书。
\end{remark}

\begin{remark}[Joukowski 压缩与有限维 toy-RH 字典的对齐]\label{rem:xi-joukowski-toy-rh-alignment}
命题 \ref{prop:xi-rh-interval-criterion} 的变量替换
\(
y=L^{s-1/2}
\)
与
\(
S=y+y^{-1}
\)
可视为 Joukowski 映射在 Dirichlet 坐标下的实现：它把单位圆 \(|y|=1\) 压缩到实区间 \([-2,2]\)，并把端点
\(
S=\pm 2
\)
对应到
\(
y=\pm 1
\)。
因此，“零点是否落在临界线 \(\Re(s)=\tfrac12\)”在 \(y\)-坐标中等价于“相应根是否落在单位圆”，而在 \(S\)-坐标中等价于“相应根是否落在 \([-2,2]\)”。
\par
该几何压缩与附录 \ref{app:real-input-40-finite-rh} 的有限维 \(\zeta\) 字典完全同型：对谱半径 \(\lambda=\rho(M)\) 的有限维核，
把 Dirichlet 变量写作 \(z=\lambda^{-s}\) 后，其临界带极点位置由归一化谱半径
\(
|\mu|/\sqrt{\lambda}
\)
控制，toy-RH（或 Ramanujan 半径界）等价于所有非平凡谱点满足 \(|\mu|\le \sqrt{\lambda}\)（命题 \ref{prop:finite-rh-triangle-dict} 与推论 \ref{cor:finite-rh-halfstrip-nopoles}）。
在“达界”情形 \(|\mu|=\sqrt{\lambda}\) 下，令
\(
y:=\mu/\sqrt{\lambda}
\)
则 \(|y|=1\) 且 \(S=y+y^{-1}\in[-2,2]\)，与本节的 \((y,S)\) 压缩坐标一致。
特别地，真实输入 \(40\) 状态核的唯一达界特征值为 \(-\sqrt{\lambda}\)（注 \ref{rem:finite-rh-40-ramanujan}），对应 \(y=-1\) 与 \(S=-2\)，从而“单频临界振荡”可被理解为该 Joukowski 压缩在端点的饱和（并见注 \ref{rem:finite-rh-40-hidden-minusA0} 对其结构来源的进一步分解）。
\end{remark}

\begin{remark}[半角提升的端点分解范式（与 Joukowski 压缩同型）]\label{rem:xi-half-angle-endpoint-paradigm}
除 \((y,S)=(L^{s-1/2},\,y+y^{-1})\) 的 Joukowski 压缩外，本文在 unitary twist 变量 \(u=e^{i\theta}\) 上还存在一个完全同型的“半角完成化”压缩：取半角变量 \(\mathfrak{r}=e^{i\theta/2}\) 使 \(u=\mathfrak{r}^2\)，并令 \(s=\mathfrak{r}+\mathfrak{r}^{-1}\in[-2,2]\)。
则幂半群 \(u\mapsto u^m\) 在完成化坐标中严格对应 \(s\mapsto C_m(s)\)（Chebyshev--Adams），并且端点 \(u=-1\)（\(s=0\)）处出现可审计的奇/偶通道喷射分裂：奇 \(m\) 通道控制线性主尺度，偶 \(m\) 通道只进入高阶喷射与有限部；同时 primitive 多项式的奇长度扇区在 \(u=-1\) 处严格消失，从而端点主奇性由偶长度扇区独占。
该端点分解的统一表述见附录 \ref{app:unit-circle-phase-gate} 的命题 \ref{prop:endpoint-decomposition-paradigm}。
\end{remark}

